\chapter*{General Introduction}
\addcontentsline{toc}{chapter}{General Introduction}
\markboth{General Introduction}{}

Mobility is one of the most basic conditions of modern economic and social life. Access to reliable transportation determines whether students reach their universities on time, whether workers can take a job in another city, and whether families can maintain regular contact across long distances. In Tunisia, however, this everyday need collides with a transportation system that has not kept pace with the growth of the population or with the geographic distribution of universities, industrial zones, and tourist regions. Public transport is often overcrowded and irregular, fuel prices have risen steadily, and inter-city travel remains expensive and time-consuming for those who do not own a vehicle.

Carpooling has emerged worldwide as a natural answer to this kind of imbalance. By matching drivers who already plan a trip with passengers who are heading in the same direction, it transforms unused seats into a shared resource, reduces the cost of mobility for both sides, and lowers the environmental impact of individual driving. In Tunisia, the demand for such a service is clearly present: thousands of trips between cities such as Tunis, Sfax, Sousse, and the smaller towns of the coast and the interior are organized every week. The supply, however, remains almost entirely informal, scattered across Facebook groups, private messages, and word-of-mouth contacts.

This informal organization carries real consequences. There is no verification of identity, no record of payment, no structured search, and no mechanism to evaluate a driver or a passenger after a trip. International platforms such as BlaBlaCar offer a more structured experience, but their adoption in Tunisia is held back by payment models that do not match local habits, by interfaces that are not localized, and by a user base that is too small to ensure frequent matches on the most requested national routes. The result is a paradoxical situation: a strong demand for shared mobility coexists with the absence of a digital tool truly designed for the country.

It is precisely this gap that the present project addresses. \textbf{ForsaDrive} is a carpooling platform conceived from the start for the Tunisian context. It is built as a unified ecosystem composed of a web application, a mobile application, and a shared backend, so that the same business rules apply regardless of the device used. The platform supports the complete lifecycle of a shared trip: account creation, identity and student verification, ride publication, search and booking, in-app communication, payment with a fifty-percent prepayment that confirms the booking, ratings, and complaint handling. It also introduces features that respond to specific local realities, including a verified student status that grants an automatic fifty-percent discount, support for cash payment of the remaining balance, and an interface designed for English, French, and Arabic users. The project was carried out within the company \textbf{ATOMIC IT}, which provided the professional environment, the technical supervision, and the methodological framework needed to take the idea from a sketch to a structured deliverable.

The remainder of this report is organized into five chapters. The first chapter presents the host organization and the general framework of the project, including a critical analysis of the current transportation situation in Tunisia and the motivations that led to ForsaDrive. The second chapter surveys existing solutions, identifies the actors of the system, formalizes the functional and non-functional requirements, and explains the agile methodology followed during the internship. The third chapter covers the conception of the platform: the general architecture, the most representative use cases, the dynamic behavior through sequence and activity diagrams, and the structural design through a class diagram and a relational schema. The fourth chapter describes the realization of the web platform, from the development environment and the technology stack to the main delivered interfaces. The fifth chapter presents the mobile application built with Flutter, the mobile-specific features it adds, and the testing strategy applied across the full platform. The report closes with a general conclusion that summarizes what was achieved, identifies the remaining limitations, and outlines the directions for future work.